% Compile with: uplatex file.tex && dvipdfmx file.dvi
%           or: latexmk -uplatex file.tex
%\documentclass[uplatex,dvipdfmx]{jsarticle}
\documentclass[dvipdfmx,a4paper,11pt]{jsarticle}
\usepackage{amsmath, amssymb}
\usepackage{tikz}

\title{scenario\_heisenberg\_two\_site}
\date{\today}

\begin{document}
\maketitle

2サイトのスピン-1/2 ハイゼンベルクモデルを扱う。結合定数 $J$（$J>0$ は反強磁性結合、$J<0$ は強磁性結合）と、両サイトで使うスピン演算子 $S$ を準備する。

J を実数パラメータとして定義して

\[
  J
\]

S をスピン演算子として定義して

\[
  S, S^{x}, S^{y}, S^{z}
\]

昇降演算子 $S^\pm=S^x\pm iS^y$ を、サイト1・2それぞれについて既存の $S^x,S^y$ の一次結合として明示的に定義する。

$Sp_1 = Sx_1 + i Sy_1$

\begin{equation}
  Sp_1 = Sx_1 + i Sy_1
  \label{eq:E1}
\end{equation}

$Sm_1 = Sx_1 - i Sy_1$

\begin{equation}
  Sm_1 = Sx_1 - i Sy_1
  \label{eq:E2}
\end{equation}

$Sp_2 = Sx_2 + i Sy_2$

\begin{equation}
  Sp_2 = Sx_2 + i Sy_2
  \label{eq:E3}
\end{equation}

$Sm_2 = Sx_2 - i Sy_2$

\begin{equation}
  Sm_2 = Sx_2 - i Sy_2
  \label{eq:E4}
\end{equation}

ハイゼンベルクハミルトニアン $H=J\mathbf S_1\cdot\mathbf S_2=JS_1^zS_2^z+\frac{J}{4}(S_1^xS_2^x+S_1^yS_2^y)$ を登録する（第2項は $S^x,S^y$ で表した $S_1^+S_2^-+S_1^-S_2^+$ に相当する）。

$H = J Sz_1 Sz_2 + J/4 (Sx_1 Sx_2 + Sy_1 Sy_2)$

\begin{equation}
  H = J Sz_1 Sz_2 + J/4 (Sx_1 Sx_2 + Sy_1 Sy_2)
  \label{eq:E5}
\end{equation}

全 $z$ 方向スピン $S_z^{\rm tot}=S_1^z+S_2^z$ を登録する。

$Sztot = Sz_1 + Sz_2$

\begin{equation}
  Sztot = Sz_1 + Sz_2
  \label{eq:E6}
\end{equation}

ハミルトニアン $H$ と全スピン $S_z^{\rm tot}$ の交換子を、それぞれ先ほど付けた名前で参照して計算する（ブラケット記法は E番号と同じ仕組みで名前付きシンボルも解決する）。ゼロになれば、全 $z$ 方向スピンは保存量である（SU(2) 対称性に由来）。

[H, Sztot]

\begin{equation}
  0
  \label{eq:E7}
\end{equation}

個々のサイトのスピンは必ずしも保存しない。サイト1の $S^z$ に対するハイゼンベルク方程式（運動方程式）を求め、非ゼロの結合項が現れることを確認する。

$Sz_1$ の Heisenberg 方程式 E5

\begin{equation}
  \frac{d Sz_1}{d t} = - \frac{1}{2} J Sx_1 Sy_2 + \frac{1}{2} J Sx_2 Sy_1
  \label{eq:E8}
\end{equation}

同様に、サイト1の $S^x$ に対するハイゼンベルク方程式を求める。

$Sx_1$ の Heisenberg 方程式 E5

\begin{equation}
  \frac{d Sx_1}{d t} = - 2 J Sy_1 Sz_2 + \frac{1}{2} J Sy_2 Sz_1
  \label{eq:E9}
\end{equation}

\end{document}